\documentclass{article}
\usepackage[a4paper, margin=1in]{geometry} % Adjust margins
\usepackage{amsmath, amssymb, amsthm}
\usepackage{tikz}
\usepackage{xcolor}
\usepackage{graphicx}
\usepackage{listings}
\usepackage{fancyhdr}
\usetikzlibrary{automata, positioning}
\usepackage{helvet} % Match font to reference
\pagestyle{fancy}
\usepackage{enumitem,xcolor}


% Ensure proper header spacing
\setlength{\headheight}{15pt} % Adjust height of the header
\setlength{\headsep}{10pt} 

% Define Header/Footer
\fancyhf{}
\fancyhead[L]{\textbf{Advanced Topics in Computer Science}} % Left header
\fancyhead[R]{Chomsky Normal Form Converter} % Right header
\cfoot{\thepage} % Page number at center bottom

% Define colors for section and subsection titles
\usepackage{titlesec}
\definecolor{sectioncolor}{rgb}{0.1, 0.3, 0.7} % Dark blue
\titleformat{\section}{\large\bfseries\color{sectioncolor}}{\thesection}{1em}{}
\titleformat{\subsection}{\bfseries\color{sectioncolor}}{\thesubsection}{1em}{}

% Define custom colors for different text types
\usepackage{tcolorbox}
\tcbuselibrary{theorems}
\newcommand{\important}[1]{\textcolor{red}{#1}}
\newtcbtheorem
  []% init options
  {definition}% name
  {Definition}% title
  {%
    colback=blue!5,
    colframe=blue!35!black,
    fonttitle=\bfseries,
  }% options
  {def}% prefix
\newcommand{\note}[1]{\textcolor{green!50!black}{#1}}

% Theorem environment setup
\newtheorem{theorem}{Theorem}
\newtheorem{lemma}{Lemma}
\newtheorem{corollary}{Corollary}

\definecolor{codegray}{gray}{0.9}
\lstdefinestyle{customc}{
  backgroundcolor=\color{codegray},
  language=C,
  basicstyle=\ttfamily\footnotesize,
  keywordstyle=\color{blue},
  commentstyle=\color{green!50!black},
  stringstyle=\color{red}
}
\begin{document}

\title{\textbf{Turing Machine Converter}}
\author{Doodlebob}
\date{}
\maketitle
\thispagestyle{fancy}

The following is an auto-explainer for Turing Machines (TM) as represented by a series of figures depicting the tape as it is moving. \\

A computation of $T = (Q, \Sigma, \Gamma, \delta, q_0, h_a, h_r)$ on input string w is a sequence of configu-
rations, where:
\begin{itemize}
\item The first configuration is $q_0w$.
\item Each configuration yields the next via a single transition. To observe what character to read,
the TM looks at the cell to the right of where the state name is.
\item $T$ is in state $q$ observing symbol $y_1$.
\end{itemize}
If the transition moves L in the leftmost cell on the tape, the TM automatically stops and rejects.

\section{Turing Machine Definition}

A Turing Machine is defined as a 7-tuple 
\[
(Q, \Sigma, \Gamma, \delta, q_0, h_a, h_r)
\]
where:
\begin{itemize}
    \item $Q$ is a finite set of states,
    \item $\Sigma$ is a finite input alphabet,
    \item $\Gamma$ is a finite tape alphabet with $\Sigma \subset \Gamma$, and $\sqcup \in \Gamma \setminus \Sigma$ is the blank symbol,
    \item $q_0 \in Q$ is the start state,
    \item $h_a, h_r \in Q$ are the accept and reject states respectively (with $h_a \neq h_r$),
    \item $\delta$ is the transition function.
\end{itemize}
\section{Transition Function}

The transition function $\delta$ is defined as follows:
\begin{align*}\delta(q0,\,0) &= (x,\,R,\,q1)\\[0.5em]\delta(q0,\,1) &= (x,\,R,\,q1)\\[0.5em]\delta(q0,\,⊔) &= (⊔,\,R,\,q2)\\[0.5em]\delta(q1,\,0) &= (0,\,R,\,q0)\\[0.5em]\delta(q1,\,1) &= (1,\,R,\,q0)\\[0.5em]\delta(q1,\,⊔) &= (⊔,\,R,\,q2)
\end{align*}

\section{Turing Machine Simulation Visualizations}

The simulation produced the following sequence of transitions. For each transition, the tape is displayed in a fixed window (so the tape does not shift), and the head (with its state label) moves along the tape.
\subsection*{Transition 1}
\begin{figure}[h]
    \centering
    \includegraphics[width=0.8\linewidth]{plots/plot1.png}
    \caption{Configuration at Transition 1}
\end{figure}
\subsection*{Transition 2}
\begin{figure}[h]
    \centering
    \includegraphics[width=0.8\linewidth]{plots/plot2.png}
    \caption{Configuration at Transition 2}
\end{figure}
\subsection*{Transition 3}
\begin{figure}[h]
    \centering
    \includegraphics[width=0.8\linewidth]{plots/plot3.png}
    \caption{Configuration at Transition 3}
\end{figure}
\subsection*{Transition 4}
\begin{figure}[h]
    \centering
    \includegraphics[width=0.8\linewidth]{plots/plot4.png}
    \caption{Configuration at Transition 4}
\end{figure}
\subsection*{Transition 5}
\begin{figure}[h]
    \centering
    \includegraphics[width=0.8\linewidth]{plots/plot5.png}
    \caption{Configuration at Transition 5}
\end{figure}
\subsection*{Transition 6}
\begin{figure}[h]
    \centering
    \includegraphics[width=0.8\linewidth]{plots/plot6.png}
    \caption{Configuration at Transition 6}
\end{figure}
\subsection*{Transition 7}
\begin{figure}[h]
    \centering
    \includegraphics[width=0.8\linewidth]{plots/plot7.png}
    \caption{Configuration at Transition 7}
\end{figure}
\end{document}